%!TEX root = hems_proj.tex
%%%%%%%%%%%%%%%%%%%%%%%%%%%%%%%%%%%%%%%%%%%%%%%%%%%%%%%%%%%%%%%%%%%%%%%
\chapter{Irodalmi Attekintes es Fogalmi Alapok}\label{ch:ALAPOK}
%%%%%%%%%%%%%%%%%%%%%%%%%%%%%%%%%%%%%%%%%%%%%%%%%%%%%%%%%%%%%%%%%%%%%%%

\begin{osszefoglal}
Áttekintjük a HEMS terület fő irányait, valamint bevezetünk néhány alapfogalmat a modell értelmezéséhez.
\end{osszefoglal}

\section{Irodalmi háttér}
A HEMS rendszerek optimalizálása tipikusan korlátozott, nemlineáris feladat, amelyben a helyi optimumok és a bizonytalan bemenetek miatt a metaheurisztikus megoldások széles körben elterjedtek \cite{Ismail2014}. A sztochasztikus termelés és a tárolás kezelési nehézségei a rendszer rugalmasságát és korlátkezelést helyezik előtérbe \cite{Fadaee2012}. A gyakorlatban ez azt jelenti, hogy a különböző algoritmusok trade-offokkal dolgoznak: gyorsaság vs. pontosság, illetve átlagos vs. legjobb megoldás.

\section{Definicio: Peak-to-Average Ratio}
\begin{ert}[Peak-to-Average Ratio -- PAR]
A PAR ertek a csucsfogyasztas es az atlagfogyasztas hanyadosa, amely a halozati terheles ingadozasat jelzi:
\begin{equation}
PAR = \frac{\max\limits_{t \in T} P_{grid}(t)}{\frac{1}{T} \sum\limits_{t=1}^{T} P_{grid}(t)}
\label{eq:par_def}
\end{equation}
Minel alacsonyabb a PAR, annal kiegyensulyozottabb a fogyasztasi profil.
\end{ert}

\begin{tet}
Ha $P_{grid}(t) \ge 0$ minden $t \in T$ idopontban, akkor $PAR \ge 1$.
\end{tet}

\begin{proof}
A definicio szerint az atlagfogyasztas nem lehet nagyobb, mint a csucsfogyasztas. Ezert
\(
\frac{1}{T} \sum_{t=1}^{T} P_{grid}(t) \le \max_{t \in T} P_{grid}(t)
\), ami alapjan a hanyados legalabb 1.
\end{proof}

\begin{pld}
Ha egy adott napon a csucsfogyasztas 4 kW, az atlag pedig 2 kW, akkor $PAR = 2$. Ez egyenetlen terhelest jelez, amelyet terheles-eltolassal csokkenteni lehet.
\end{pld}

\begin{meg}
A PAR csokkentese nem csak koltsegcsokkentessel jar, hanem a halozati csucsok merseklesevel is, ami a szolgaltato szamara is elonyos.
\end{meg}
